% =====================================================================
% DRAFT REPLACEMENT for the real-vehicle results, discussion and limitations.
%
% NOT INCLUDED IN THE BUILD. Written as a separate file because the fixer agent
% was editing the manuscript concurrently; merge into elsarticle-template-num.tex
% once that build is verified.
%
% Every number below traces to reports/matrix/folds_table_symmetric.md, produced by
%   <py311> scripts/eval_folds_table.py --selected
% and to reports/matrix/selected_checkpoints.json. Curves for the figure are in
% reports/matrix/folds_table_symmetric_curves.npz; the figure itself is
% figures/fig_main_cv.pdf from scripts/fig_main_cv.py.
%
% WHAT CHANGED RELATIVE TO THE VERSION ON DISK, AND WHY
%   * Every value is new. The old table reported ONE checkpoint per model on ONE
%     split; controlled retrains showed both our row and the baseline rows were
%     favourable single draws, so nothing from it survives.
%   * The headline claim changes from "1.4x lower error" to "far less sensitive to
%     which episodes are held out", because that is what the 5-fold data actually
%     shows and it does not depend on a single lucky fit.
%   * The delta-noise subsection inverts: the paper called the non-monotonicity
%     single-run variance; across five folds the trend is monotone and reproducible.
%     But a matched-magnitude Gaussian control reproduces the gain, so this is
%     ordinary regularisation and must NOT be sold as a property of the
%     conference's weak-supervision noise model.
%   * SINDYc is now measured per fold instead of carried over from the old split.
%
% STILL MISSING -- do not merge without resolving:
%   [ ] Table 2 (kappa-source ablation) is still single-split, from the same
%       un-reproducible dyn_k16 checkpoint family as the old main table. Either
%       re-run the ablation per fold (needs finetune/frozen/vqformer encoders for
%       each of the 5 folds, ~15 jobs) or mark that table explicitly as
%       single-split and not comparable to Table 1.
%   [ ] The factorized-state dimension (R^6) contradicts the dynamics equations and
%       the implementation. Method section fix, tracked separately.
% =====================================================================

\subsection{Real-vehicle results}\label{sec:donkey_results}

All numbers in this subsection come from five-fold cross-validation over the
$71$ recorded episodes. Each fold holds out one fifth of the episodes
($14$--$15$ episodes, $301$--$647$ evaluation windows), and every model is
trained on that fold's training episodes only. Because the folds hold out
different episodes, window counts differ between them; we therefore average
within a fold first and report the mean and half-range over the five fold-level
means. Comparisons \emph{within} a fold are paired --- every model is rolled out
from the same initial states over the same windows in the same order --- and all
paired statistics below are computed that way. Comparisons across folds are not
paired and we do not treat them as such.

Cross-validation restores the protocol of the conference version, which evaluated
its variants under five-fold cross-validation; the single-split evaluation used in
an earlier draft of this article was a step back from it, and the paired
within-fold tests below go beyond what either reported.

\paragraph{Checkpoint selection}
Our dynamics keeps the epoch with the lowest $E_{xy}(100)$, the quantity we
report. Selecting the baselines by their composite validation loss instead, as
an earlier version of this evaluation did, penalises them: re-selecting each
baseline by the reported metric improves DVBF from \SI{1.287}{m} to
\SI{0.837}{m} and GOKU-net from \SI{0.815}{m} to \SI{0.725}{m}. We therefore
apply the same rule to every model, choosing each baseline's epoch from
saved snapshots by $E_{xy}(100)$. This makes selection symmetric; it does not
make it unbiased, since both sides then select on the evaluation windows. It also
differs from the conference version, which stopped every model early on validation
loss; that rule is the more conservative of the two and we note the change
explicitly rather than let the two protocols be read as one.

% ---------------------------------------------------------------- main table
\begin{table}[t]
\centering
\caption{Position error $E_{xy}(k)$ in metres on the real vehicle, mean $\pm$
half-range over five folds. Lower is better. Every model is selected by
$E_{xy}(100)$ and rolled out for $100$ steps from the ground-truth state. The
baselines lead by a wide margin at short horizons; the ordering inverts before
step $100$. The rightmost column also shows the sensitivity of each model to
which episodes are held out.}
\label{tab:donkey_main}
\begin{tabular}{lccc}
\toprule
Model & $E_{xy}(25)$ & $E_{xy}(50)$ & $E_{xy}(100)$ \\
\midrule
Vid2Param~\cite{asenov2019vid2param}  & $\mathbf{0.036 \pm 0.010}$ & $\mathbf{0.107 \pm 0.009}$ & $0.703 \pm 0.192$ \\
GOKU-net~\cite{linial_generative_2021} & $0.046 \pm 0.015$ & $0.131 \pm 0.015$ & $0.725 \pm 0.122$ \\
DVBF~\cite{karl2016deep}              & $0.061 \pm 0.011$ & $0.173 \pm 0.038$ & $0.837 \pm 0.160$ \\
SINDYc~\cite{BRUNTON2016710}          & \multicolumn{3}{c}{diverges in all $2483$ windows of all five folds} \\
\midrule
\textbf{PIWM-Frenet (ours)}           & $0.141 \pm 0.011$ & $0.245 \pm 0.025$ & $\mathbf{0.463 \pm 0.046}$ \\
\textit{PIWM-Frenet (map-free)}       & $0.146 \pm 0.013$ & $0.258 \pm 0.031$ & $0.516 \pm 0.068$ \\
\bottomrule
\end{tabular}
\end{table}

\paragraph{Long-horizon prediction}
Table~\ref{tab:donkey_main} and Figure~\ref{fig:donkey_main} show the same
crossover the conference version reported, now with the spread across folds made
visible. Up to roughly step $65$ the baselines are ahead, and by a large factor:
at step $25$ Vid2Param's error is a quarter of ours. This is not noise --- its
half-range at that horizon is \SI{0.010}{m}. Beyond the crossover the ordering
inverts and does not invert back. At step $100$ our model reaches
\SI{0.463}{m} against \SI{0.703}{m} for the strongest baseline. Within each
fold, where the comparison is paired, our model is ahead at step $100$ in
\emph{all five} folds against every baseline: by \SI{0.241}{m} on average
against Vid2Param, \SI{0.262}{m} against GOKU-net and \SI{0.374}{m} against
DVBF. SINDYc diverges in every window of every fold; we report it as divergent
rather than quoting a value, since past the point of divergence the number
reflects only where the rollout was clamped.

\paragraph{Sensitivity to the held-out episodes}
The more informative quantity is the last column's spread. Changing which
episodes are held out moves our error by $\pm\SI{0.046}{m}$, against
$\pm\SI{0.192}{m}$ for Vid2Param, $\pm\SI{0.122}{m}$ for GOKU-net and
$\pm\SI{0.160}{m}$ for DVBF --- a factor of three to four. The baselines' error
at long horizons depends substantially on which part of the track they were
trained on; ours much less so. We regard this as the more durable claim of the
two, because it does not rest on any single trained model: the same conclusion
holds in every fold, whereas a single-split comparison of best checkpoints can
be, and in our earlier evaluation was, dominated by which draw each side
happened to get.

% ------------------------------------------------------------- delta noise
\paragraph{Weak supervision with noisy labels}
Table~\ref{tab:delta} repeats the evaluation with the conference weak-supervision
noise applied to the training labels at $\delta \in \{0, 5\%, 10\%\}$. The two
model families respond in opposite directions. Every baseline degrades as the
labels get noisier. Our model improves, monotonically, from \SI{0.463}{m} at
$\delta = 0$ to \SI{0.319}{m} at $\delta = 10\%$ --- a $31\%$ reduction, while
Vid2Param worsens by $33\%$ over the same range. The trend is consistent across
all five folds and the intervals do not overlap.

\begin{table}[t]
\centering
\caption{$E_{xy}(100)$ in metres under weak-supervision label noise, mean $\pm$
half-range over five folds. The structured model improves as the labels degrade;
every baseline gets worse.}
\label{tab:delta}
\begin{tabular}{lccc}
\toprule
Model & $\delta = 0$ & $\delta = 5\%$ & $\delta = 10\%$ \\
\midrule
\textbf{PIWM-Frenet (ours)}     & $0.463 \pm 0.046$ & $0.452 \pm 0.047$ & $\mathbf{0.319 \pm 0.049}$ \\
\textit{PIWM-Frenet (map-free)} & $0.516 \pm 0.068$ & $0.506 \pm 0.066$ & $0.390 \pm 0.040$ \\
Vid2Param                       & $0.703 \pm 0.192$ & $0.919 \pm 0.252$ & $0.938 \pm 0.280$ \\
GOKU-net                        & $0.725 \pm 0.122$ & $0.896 \pm 0.199$ & $0.864 \pm 0.185$ \\
DVBF                            & $0.837 \pm 0.160$ & $1.067 \pm 0.240$ & $1.075 \pm 0.233$ \\
\bottomrule
\end{tabular}
\end{table}

\paragraph{The gain is regularisation, not the noise model}
It would be convenient to read Table~\ref{tab:delta} as evidence that the
particular weak-supervision noise of the conference formulation benefits a
structured model. It does not. Replacing that biased-uniform noise with a
zero-mean Gaussian of the same per-dimension magnitude, holding everything else
fixed, gives \SI{0.313 \pm 0.043}{m} against \SI{0.319 \pm 0.049}{m} --- a
paired difference of \SI{-0.005}{m} whose sign is not even consistent across
folds. Label noise at this magnitude regularises the structured model, and any
zero-mean noise of that magnitude does it. The finding is the \emph{asymmetry}
between the model families, not the form of the noise. We note that the
baselines' selected epochs move sharply earlier as $\delta$ grows --- several
select the fifth epoch at $\delta = 10\%$ --- which is consistent with their
fitting the label noise almost immediately.

% ------------------------------------------------------------- map-free
\paragraph{The cost of dropping the known centerline}
The rows labelled map-free obtain the vehicle's position without consulting the
track map at all, reading the road geometry from the perception's shape output
instead of indexing a surveyed centerline. Doing so costs \SI{0.053}{m} at step
$100$, paired within folds and consistent in sign across all five
($+0.049$, $+0.022$, $+0.069$, $+0.071$, $+0.056$). This is the version of the
model that could be deployed on a track that has not been measured, and it
remains ahead of every baseline at step $100$.

% =====================================================================
% DISCUSSION -- replacement for the corresponding paragraphs
% =====================================================================

\paragraph{What the crossover means}
A model that is four times worse at step $25$ and a third better at step $100$
is not uniformly better, and we do not claim it is. The baselines fit the local
motion more tightly; what they do not do is stay bounded. Our error grows and
then flattens, because the road is read once from the camera and re-indexed by
arc length rather than extrapolated, so there is no compounding road prediction
to diverge. For a controller that needs the next few tenths of a second, the
baselines are the better choice on this data. For anything that must reason a
few seconds ahead, they are not.

\paragraph{Structure rather than information}
Both model families see the same images and the same actions. The map-free
result sharpens this: removing the surveyed centerline entirely costs
\SI{0.053}{m}, which leaves the long-horizon gap largely intact. What
distinguishes the two families is therefore not what they are given but how the
road enters the state --- as a quantity that is advected geometrically, or as
one that must be predicted forward and can drift.

% =====================================================================
% LIMITATIONS -- replacement
% =====================================================================

% NOTE ON MERGE: the existing Limitations list has a fifth item stating that each
% configuration is trained once on a single split and that a cross-validated
% repetition is needed. DELETE it -- that repetition is what this section now
% reports. Leaving it in would contradict the tables directly.
\paragraph{Limitations}
Four limitations bound what these results support.

First, the evidence is from a single vehicle on a single track. The
cross-validation varies which episodes are held out, not which track or which
platform, so it measures sensitivity to the driving data and not to the
environment. Whether the advantage transfers to other road geometries is
untested here.

Second, every model's epoch is selected on the same held-out windows the table
reports. We apply the rule symmetrically, which removes the asymmetry present in
our earlier evaluation, but symmetric selection is not unbiased selection and
the absolute values should be read with that in mind.

Third, only one baseline receives images. Vid2Param infers a parameter vector
from the image window because that is the whole of the method; DVBF and GOKU-net
are run as transition models on the shared state. Restoring GOKU-net's
observation pathway in the obvious way does not produce an independent
comparison --- with a parameter head of matched width it becomes architecturally
identical to Vid2Param and, trained under the same seed, reproduces it exactly.
A genuinely independent observation-conditioned GOKU-net would require
implementing its latent-ODE formulation, which we did not do.

Fourth, the improvement under label noise is ordinary regularisation, as the
Gaussian control shows. We report the asymmetry between model families, not a
property of the weak-supervision formulation, and the mechanism behind that
asymmetry is not established by these experiments.
